\subsection{Adjoint sensitivity vs.\ direct autograd}
\label{sec:track-a}

Our own implementation documentation adopted direct autograd on the grounds
that the adjoint method's $O(1)$ memory advantage
(Section~\ref{sec:gradients}) does not matter for our short trajectories
($N_t\in[36,101]$). We test that claim directly on an RTX~3060. The
integrator is held fixed (classical RK4 at matched step size on both paths),
and a forward-trajectory check confirms the two implementations agree to
$\sim10^{-7}$ relative error, so only the differentiation method varies.

\begin{expectbox}
\expectrow{Question}{Does adjoint sensitivity's constant memory cost make it the better choice at our trajectory lengths, as our own documentation assumes?}
\expectrow{Expected}{Direct autograd should be cheaper in memory at these short lengths, with adjoint's advantage only appearing once trajectories grow much longer.}
\expectrow{Observed}{Adjoint already uses less memory at every length tested. The deciding cost is wall-clock, not memory: adjoint is $1.9$--$2.1\times$ slower.}
\end{expectbox}

\begin{table*}[t]
\centering
\small
\renewcommand{\arraystretch}{1.15}
\caption{Peak VRAM and wall-clock, direct autograd vs.\ adjoint sensitivity.}
\label{tab:adjoint-profiling}
\begin{tabularx}{\textwidth}{@{}X C{2.8cm} L{3.4cm} C{2.8cm} C{3.2cm}@{}}
\toprule
\hdrrow \thh{System} & \thh{Trajectory length $N_t$} & \thh{Method} & \thh{Peak VRAM (MB)} & \thh{Seconds per 1000 epochs}\\
\midrule
Lotka-Volterra & 36 & Direct autograd & 18.61 & 317.5\\
Lotka-Volterra & 36 & Adjoint & 17.08 & 672.9\\
SIR & 101 & Direct autograd & 21.52 & 984.8\\
SIR & 101 & Adjoint & 17.12 & 1871.6\\
\bottomrule
\end{tabularx}
\end{table*}

\begin{figure}[tb]
\centering
\includegraphics[width=\columnwidth]{figures/track_a/memory_vs_trajectory_length.png}
\caption{Peak memory vs.\ trajectory length $N_t\in\{36,\dots,3201\}$: direct
autograd grows linearly ($\approx0.045$~MB per output point); adjoint stays
flat (0.7\% increase over an $89\times$ length increase).}
\label{fig:adjoint-memory}
\end{figure}

\paragraph{Finding.} The measured memory data contradicts the premise behind
the original choice, but still supports its conclusion. The $O(N_t)$ vs.\
$O(1)$ scaling in Equation~\ref{eq:adjoint} predicts a crossover length below
which direct autograd should still win on raw memory, because it avoids the
adjoint method's own bookkeeping overhead. The data shows no such crossover
inside the tested range. Adjoint uses \emph{less} memory than direct autograd
at every tested length, including the shortest ($N_t=36$: $17.08$ vs.\
$18.61$~MB). As Figure~\ref{fig:adjoint-memory} shows, direct autograd's
per-step storage cost is already large enough at our scale to outweigh
whatever constant overhead the adjoint bookkeeping carries.
Extrapolating the measured linear slope of direct autograd's memory growth,
$N_t$ would need to reach $\approx225{,}000$ before it threatens a 12~GB
card. For us the deciding factor is wall-clock: solving the continuous
adjoint ODE backward in time costs a consistent $1.9$--$2.1\times$ overhead
at both tested lengths. That turns a 53-minute Lotka-Volterra run into
$\sim$112~minutes, for a memory saving that never matters in practice (peak
usage is $0.17\%$ of card capacity). Gradient relative error is
$2.85\times10^{-7}$ and $4.88\times10^{-7}$ respectively, well inside a
$5\times10^{-3}$ tolerance. Both methods agree numerically, so the choice
reduces to a resource trade-off, and direct autograd is the faster option
for every dataset in our scope ($N_t\le201$).
